% ============================================================
%  CHAPTER 10 — CONCLUSION AND FUTURE WORK
% ============================================================
\chapter{Conclusion and Future Work}

\section{Conclusion}

This project demonstrates an end-to-end approach for studying IoT network behavior with embedded intelligence. Aegis-IoT integrates agent-based simulation, predictive modeling, anomaly scoring, adaptive control, and multi-surface visualization into a single workflow that supports repeatable experiments.

\section{Key Outcomes}

\begin{itemize}
    \item A structured AnyLogic simulation of IoT packet flows with dataset-calibrated telemetry.
    \item Inline inference for latency prediction, anomaly scoring, and short-horizon forecasting.
    \item A tabular Q-learning policy to adapt behavior under varying congestion levels.
    \item A federated learning prototype demonstrating parameter aggregation without raw data exchange.
    \item Desktop and web dashboards for real-time monitoring and offline analysis.
\end{itemize}

\section{Limitations}

\begin{itemize}
    \item The experiments are simulation-based and use dataset-derived distributions rather than a live MQTT broker.
    \item The energy and stress-traffic models are simplified proxies intended for controlled evaluation.
    \item Tabular Q-learning scales to small discrete spaces; richer state representations would require function approximation.
\end{itemize}

\subsection{Future Work}

\begin{itemize}
    \item Integrate live MQTT traffic ingestion to validate models under real network conditions.
    \item Extend the controller from tabular Q-learning to deep RL for continuous or high-dimensional state.
    \item Improve forecasting with sequence models and evaluate robustness to distribution shift.
    \item Add privacy mechanisms (e.g., differential privacy) to the federated learning prototype.
\end{itemize}
